
\vspace{-0.08in}
\section{Discussion and Limitations}
\label{append:discussion}

\noindent \textbf{Can the defined security properties and four attacks guarantee discovery of all vulnerabilities?} 
\lm aims to reduce reliance on slow, manual expert review by enabling systematic and scalable detection of standard-level defects in the extracted function chains. It does not guarantee the discovery of all vulnerabilities for two main reasons.
First, security verification is inherently scoped. Any finite set of security properties may miss emerging or unmodeled failure modes. We can only claim security with respect to the predefined properties and assumptions of the threat model~\cite{chen2023sherlock,al2024hermes}. Second, while not all SCA nodes correspond to explicit message exchanges, some nodes represent event or condition-driven transitions such as local policy decisions. In this case, \textit{SecOracle} still applies the four attacks to such nodes by exploiting their triggered events or conditions. If an SCA node is purely deterministic without any inputs, such as entering one state after another state, this transition is non-attackable under our settings. Notably, the majority of reported vulnerabilities in RRC and NAS layers included in Table~\ref{tab:vulnerabilities} stem from message exchanges and thus fall within the scope of our attack model.   

\noindent  \textbf{Out-of-scope vulnerabilities.}
\lm currently targets standard design flaws that can be exposed through adversarial manipulation of message or event flows, including dropping, modifying, injecting, and replaying messages. This scope covers availability, authentication, authorization, integrity, confidentiality, and state-consistency issues in control-plane procedures. We admit that several important vulnerability classes are outside the current model. For example, \lm does not directly analyze cryptographic weaknesses, side channels, physical-layer or social-engineering attacks, and vendor deviations absent from specifications, and may miss deployment-specific vulnerabilities involving carrier policies or undocumented hardware.





\noindent\textbf{Specification evolution and cross-release analysis.} 

3GPP specifications continuously evolve across releases. Previously identified security issues may persist, be partially mitigated, or be eliminated by later revisions. These ongoing efforts are essential for strengthening the security of cellular systems. At the same time, the scale and complexity of the standards make it challenging to identify and eliminate every security-relevant design weakness through traditional methods. Particularly, we examined the specification text corresponding to our newly reported vulnerabilities in the latest available versions at the time of submission and found that all 7 vulnerabilities remain present, as shown in Table~\ref{tab:new_vulnerabilities}. This observation motivates the automated cross-release analysis. In future work, we plan to extend \lm to compare SCA nodes and function chains across versions, identify security-relevant changes, and determine whether previously identified weaknesses are mitigated, persist, or are reintroduced as specifications evolve.






\noindent\textbf{Ethical Consideration and Responsible Disclosure.}
Our research follows established ethical standards for security research. Identified vulnerabilities in cellular specifications are handled throu-\ gh responsible disclosure.
We work with stakeholders, including standardization organizations, operators, and device manufacturers, to report design defects. Their acknowledgments will be posted on our project website.
Our objective is not to expose flaws for exploitation, but rather to advance an efficient methodology for enabling 3GPP, operators, and vendors to rigorously review specifications prior to implementation and public deployment of services. 